\documentclass[../Main.tex]{subfiles}
\begin{document}

\section{Motivation}
In the era of social networking, people have an insatiable need for constant entertainment. Young adults, in particular, are the first generation growing up with social networking sites (SNS) like Facebook and games (SNG) like World of Warcraft. However, the extent of this is still unclear.

In the later years, there have been conducted multiple studies to find out what consequences modern technology has on young adults. Many of these studies examine the direct effects of SNS and SNG. Studies like the ones performed by \citeauthor{pantic_online_2014} in \citeyear{pantic_online_2014} and \citeauthor{mentzoni_problematic_2011} in \citeyear{mentzoni_problematic_2011}, concluded that SNS and SNG in fact do correlate to users’ mental and physical health. 

Alternative games can be defined as ... Many existing alternative games use gamification and other related concepts to motivate and engage the users in their responsibilities, such as edutainment games, exergames, and well-being games. The goal of gamification is to support value creation for the user so that he or she can get a sense of motivation, achievement, and mastery \cite{hamari_social_2013}. 

There exist multiple success stories of alternative games that have sparked from the idea of helping people to educate themselves and improve their physical and mental well-being. One such example is Strava. This application can be seen as both a SNS and SNG, because it uses gamification to engage and motivate people to work out. Users are rewarded with positive reinforcement in form social interaction in the form of likes, comments, and from achieving good results and breaking records.

For us personally, it is a great motivation to produce something that can be socially beneficial. This is also motivated by the fact that we ourselves would like to create something that we can benefit from ourselves, whether that is social benefits or health benefits.
\clearpage

\section{Project and context}
This work is the pre-project for our Master thesis where the task given to us was as follows:

\blockquote{
\textbf{[Game technology] Alternative games}

In this project, the goal is to prototype an innovative game(s) and test this game(s) on users. The innovation can be in the type of gameplay the game provides, how it combines various game genres, what technology is used to control or play the game, how the social interaction between players is supported, the purpose of the game, etc.

The first phase of the project will include a study of innovative games, game technologies, and concepts, the second phase will be to implement one or more prototypes, and the third phase will be to evaluate the concept through playtesting on real users.}

The project consists of exploring various concepts and ideas in the space of alternative games. Our work will contribute to the alternative games field of game technology. Different concepts will be explored, where the most promising one will be further evaluated and tested during the Master thesis period.

\section{Research goal, questions, and methods}
In this project, the Goal, Question, Metrics (GQM) approach \cite{basili_software_1992} is used to define the software metrics. This methodology has a conceptual, an operational and a quantitative level. At the conceptual level we define a goal for what the project should achieve. At the operational level we define a set of research questions that are associated with the goal. These questions can, if answered, indicate whether we are accomplishing our goal. The quantitative level we define metrics that help answer the research questions.


\subsection{Research goal}
\textit{The goal of this project is to develop one or more game concepts for an alternative game that has a socially beneficial effect, and that can be developed, explored, and tested.}

\subsection{Research questions}
\textbf{RQ1}: \textit{What is an alternative game?} \\
This research question investigates what alternative games are, and how they differ from other games.

\textbf{RQ2}: \textit{How can alternative games be categorized and conceptualized?} \\
This research question will examine how existing alternative games can be categorized using a structured approach. This research will highlight existing alternative games, and point to what exactly makes them alternative. This will also investigate if a structured approach will help the process of developing new concepts of alternative games.

\textbf{RQ3}: \textit{What makes a game engaging?} \\
This research question will investigate theories and concepts used in game development that makes games engaging. This research will explore concepts that can enhance the players’ enjoyment, motivation, and engagement.

\textbf{RQ4}: \textit{ Which game concept(s) can be developed with existing technology that has a socially beneficial purpose? How do this/these contribute to the study of alternative games?} \\
This research question will consider game concepts that can be developed based on the knowledge gained from literature and technology reviews. The concepts that are developed, should comply with the definition of alternative games and  have some social beneficial purpose. It is also of interest to study how these concepts can contribute to  the study of alternative games though our interpretation of alternative games. The most promising concepts will be further conceptualized and developed(at a later stage). 



\subsection{Research methods}
The research questions will be answered through an analysis of relevant literature. These areas include existing alternative games, game development, game flow, game platforms, and available game technology. \citeauthor{aveyard_doing_2014}  stated in \citeyear{aveyard_doing_2014} that a literature review is a “comprehensive study and interpretation of literature that relates to a particular topic”. The insight from this literature study on existing games, relevant technologies, and concepts of game development will be summarized in the chapters of part \ref{part:prestudy} - Prestudy (ref/label). The chapters will give both the reader and the writers the background information needed to do further research in this field. Through the review of the literature, we will seek to answer the research questions RQ1-4. This will give insight into what alternative games are, what alternative games already exist, and how they contribute to the field. We will also review relevant game development theories on how to create fun and engaging games. 

The research field of alternative games covers many types of games, and the pool of existing examples are very diverse. Each game is different from the other. The insight acquired from studying existing games in RQ2 and game theory in RQ3, will then be used to develop a model for categorizing alternative games. The structured approach of modelling will be used for finding synergies between the games. 

Furthermore, the model will then be used to test RQ3 when developing new concepts for alternative games. The most promising concepts will be developed in detail before one final concept is prototyped.



\section{Readers guide}
\subsection{Part \ref{part:introduction} - Introduction}
Introduced the project, context, and our motivation for choosing this task. In addition to this,  research goal, questions, and methods were presented.

\subsection{Part \ref{part:prestudy} - Prestudy}
Insert short description and readers guide of part \ref{part:prestudy}.

\subsection{Part \ref{part:contrubution} - Our contribution}
Insert short description and readers guide of part \ref{part:contrubution}.
Our contribution to the field of alternative game will be presented in this project, and will be useful as a pre-requirement to build upon in future projects.

\subsection{Part \ref{part:discussion} - Discussion \& conclusion}
Insert short description and readers guide of part \ref{part:discussion}.

\biblio % Needed for referencing to working when compiling individual subfiles - Do not remove
\end{document}